\chapter{The model and the estimators}\label{ch:model}

This chapter summarizes the exponential random graph model, the estimators that \nwergm{} offers, and the choices that govern how they behave: the sampler, the stopping rules, reproducibility, missing dyads and constraints. It follows \ergm{}; the names of the corresponding arguments of R's \texttt{ergm()} and \texttt{control.ergm()} are given in parentheses. The syntax of the options is in Chapter~\ref{ch:options}.

\section{The exponential random graph model}\label{sec:ergmdef}

Let $Y$ be the random adjacency matrix of a network with $n$ nodes, with $Y_{ij}=1$ if there is a tie from $i$ to $j$ (in an undirected network $Y_{ij}=Y_{ji}$; in a two-mode network the ties run between the two sets of nodes only). An ERGM specifies the probability of a network $y$ as
\begin{equation}\label{eq:ergm}
  P_\theta(Y=y)=\frac{\exp\{\eta(\theta)^{\top}g(y)\}}{\kappa(\theta)},\qquad
  \kappa(\theta)=\sum_{y'\in\mathcal{Y}}\exp\{\eta(\theta)^{\top}g(y')\},
\end{equation}
where $g(y)$ is a vector of network statistics (the \emph{terms} of the model), $\eta(\theta)$ the vector of their coefficients as a function of the parameters $\theta$, and $\mathcal{Y}$ the set of networks on which the model is defined (all binary networks on the given nodes, unless constraints restrict it, Section~\ref{sec:constraints}). In an ordinary model $\eta(\theta)=\theta$. In a \emph{curved} model $\eta$ is a nonlinear function of $\theta$: for the geometrically weighted terms with an estimated decay, a term $\sum_k w_k(\alpha)\,g_k(y)$ has weights $w_k(\alpha)=e^{\alpha}\{1-(1-e^{-\alpha})^{k}\}$ and the model has a weight and a decay parameter (Hunter and Handcock~\cite{hunter2006}, Hunter~\cite{hunter2007}). The constant $\kappa(\theta)$ sums over all $2^{n(n-1)}$ (directed) networks and cannot be computed except for very small networks, which is what makes estimation hard.

\paragraph{Reading a coefficient.} The model implies, for every dyad, the log odds of a tie given all other dyads,
\begin{equation}\label{eq:condlogit}
  \operatorname{logit}P_\theta\bigl(Y_{ij}=1\mid y_{-ij}\bigr)=\eta(\theta)^{\top}\delta_{ij}(y),\qquad
  \delta_{ij}(y)=g(y_{ij}^{+})-g(y_{ij}^{-}),
\end{equation}
where $y_{ij}^{+}$ and $y_{ij}^{-}$ are the network with the tie set to present and to absent, and $\delta_{ij}$ is the \emph{change statistic} of the dyad. A coefficient is therefore the change in the conditional log odds of a tie when the corresponding statistic increases by one, holding the rest of the network fixed. For the term \texttt{edges} the change statistic is always 1, so that its coefficient is the log odds of a tie in a network with no other effect. For \texttt{mutual} the change statistic is 1 if the reverse tie is present and 0 otherwise, so that its coefficient is the additional log odds of a tie that reciprocates an existing one. Coefficients of dyad-dependent terms are conditional on the rest of the network and cannot be read as effects of a single covariate in isolation.

\paragraph{Dyad independence.} If every change statistic depends only on the covariates of the two nodes of the dyad (as for \texttt{edges}, \texttt{nodecov}, \texttt{nodematch}, \texttt{edgecov}), the dyads are independent, the model is a logistic regression for the tie indicators, and $\kappa(\theta)$ factorizes. Any term that depends on other ties (\texttt{mutual}, \texttt{triangle}, \texttt{gwesp}, degree terms, \ldots) makes the model dyad-dependent. The chapters on the terms record which is which.

\section{Estimators}\label{sec:estimators}

\nwergm{} offers three estimators, chosen with \opt{estimate()} (R: \texttt{estimate=}) and, for the Monte Carlo estimator, \opt{mainmethod()} (R: \texttt{main.method=}). If \opt{estimate()} is not given, the choice is made by the model: a dyad-independent model is fitted by the pseudolikelihood, which is then the exact maximum likelihood, and all other models by Monte Carlo maximum likelihood. The older option \opt{method(mple|mcmle|cd)} is a synonym.

\subsection{Maximum pseudolikelihood}\label{sec:mple}

The pseudolikelihood replaces the likelihood by the product of the conditional probabilities~\eqref{eq:condlogit} over all dyads,
\[
  \mathrm{PL}(\theta)=\prod_{i,j}P_\theta\bigl(Y_{ij}=y_{ij}\mid y_{-ij}\bigr),
\]
which is the likelihood of a logistic regression of the tie indicators on the change statistics, evaluated on the observed network. \nwergm{} builds the matrix of change statistics (in the plugin when it can) and maximizes the pseudolikelihood by Newton's method with a tight convergence tolerance, so that the coefficients agree with R's to about $10^{-6}$ or better (Chapter~\ref{ch:validation}). Perfectly collinear columns cannot be estimated and are reported as omitted. For a dyad-independent model the estimate and its standard errors are those of the maximum likelihood estimator. For a dyad-dependent model the standard errors of the pseudolikelihood are not valid and the estimate is mostly useful as a starting value for the Monte Carlo estimator; \nwergm{} says so in a note under the table. Missing dyads are left out of the product (Section~\ref{sec:missing}).

The variants (R's \texttt{control.ergm()} arguments):
\begin{description}[style=nextline,nosep]
  \item[\opt{mpletype(glm|penalized)} (\texttt{MPLE.type})] \texttt{penalized} is the Firth (Jeffreys-prior) penalized pseudolikelihood, which is finite when the ordinary estimate does not exist because of separation. It cannot be combined with \opt{offset()}.
  \item[\opt{mplecheck(on|off)} (\texttt{MPLE.check})] Tests whether the estimate exists (a linear program for separation, solved by a simplex routine). If it does not, \nwergm{} warns \texttt{The MPLE does not exist!} and, with \opt{estimate(mple)}, reports the last iterate of \texttt{glm()}, as R does (huge coefficients, standard errors in the thousands). With \opt{mplecheck(off)} the iterate is reported without the test.
  \item[\opt{mplecovariance(invhess|godambe|bootstrap)} (\texttt{MPLE.covariance.method})] Replaces the inverse-Hessian covariance of a dyad-dependent pseudolikelihood fit by the Godambe sandwich or by the covariance of 500 pseudolikelihood fits to networks simulated at the estimate (burn-in 1024, interval 1024).
\end{description}

\subsection{Monte Carlo maximum likelihood}\label{sec:mcmle}

Monte Carlo maximum likelihood (MCMLE) follows Geyer and Thompson~\cite{geyer1992} and, for the step length and the stopping rules, Hummel, Hunter and Handcock~\cite{hummel2012}. For a fixed starting value $\eta_0$, the log-likelihood ratio is
\[
  \ell(\eta)-\ell(\eta_0)=(\eta-\eta_0)^{\top}g(y_{\mathrm{obs}})-\log E_{\eta_0}\exp\{(\eta-\eta_0)^{\top}g(Y)\},
\]
and the expectation is replaced by an average over a sample of networks simulated from the model at $\eta_0$. Each outer iteration therefore (1) draws a Markov chain sample of networks at the current coefficients, (2) takes a step towards the maximiser of the approximate log-likelihood ratio, with a step length that keeps the step inside the region where the approximation can be trusted (the Hummel step length), and (3) tests whether the sample statistics are already consistent with the observed statistics. The starting values are the pseudolikelihood estimates (or the values of \opt{init()}). On the models of this manual \nwergm{} needed one to four iterations (eight for the Lazega model with the shared-partner term in Section~\ref{sec:comparing}). The number of iterations is stored in \texttt{e(mcmle\_iterations)}, and the limit is \opt{mcmleiterations()} (R: \texttt{MCMLE.maxit}, default 60); a fit that reaches the limit is reported with a warning and \texttt{e(converged)}${}=0$.

\paragraph{The sampler.} The networks are simulated by a Metropolis--Hastings chain that proposes to toggle one dyad at a time. The default proposal (\opt{proposal(tnt)}, R's \texttt{MCMC.prop.weights="TNT"}) is the \emph{tie/no-tie} proposal of Morris, Handcock and Hunter~\cite{morris2008}: with probability one half it proposes to remove a randomly chosen existing tie, and otherwise to add a randomly chosen non-tie, with the acceptance probability corrected for the unequal chance of proposing a given toggle. Because most networks are sparse, this mixes much better than the uniform proposal (\opt{proposal(uniform)}, or \opt{proposal(random)}), which picks a dyad uniformly. In a fixed-size sample (\opt{mcmleeffectivesize(0)}, and the fits by contrastive divergence, stochastic approximation and with \opt{obsconstraints()}) the chain runs for \opt{mcmcburnin()} proposals (default 16384), then records a network every \opt{mcmcinterval()} proposals (default 1024) until \opt{mcmcsamplesize()} (default 3000) networks are recorded. In the default adaptive sample of Monte Carlo maximum likelihood, described next, the number of draws is not fixed. The interval and the burn-in default to those of \ergm{}'s \texttt{control.ergm()}. Constraints change the proposal; see Section~\ref{sec:constraints}.

\paragraph{Adaptive sample size and interval.}\label{sec:adaptive} Draws of a Markov chain are autocorrelated, so that a sample of a given size carries less information than the same number of independent draws. \nwergm{} samples as \ergm{}'s \texttt{ergm()} does (R: \texttt{MCMLE.effectiveSize}, \texttt{MCMLE.samplesize}), with a sample size that adapts to the autocorrelation of the chain and not a fixed number of draws. Each Monte Carlo iteration (1) draws a base sample whose size is \opt{mcmcsamplesize()}, by default the larger of 256 and 32 times the number of coefficients (R: \texttt{MCMLE.samplesize.min}, \texttt{.per\_theta}), at the current interval; (2) estimates from the sample itself how many initial draws to discard as burn-in (a trend test on the estimating equations, as \ergm{} does); (3) estimates the effective sample size (ESS) of the estimating equations of the draws that remain; and (4) if it is below the target \opt{mcmleeffectivesize()} (default 64), draws more, in at most 16 runs whose lengths are predicted from the shortfall, and thins the whole sample by two, doubling the interval, whenever it has doubled in size. The first iteration starts at \opt{mcmcinterval()} (1024); every later iteration starts from half the interval that the previous sample ended with (R: \texttt{MCMLE.effectiveSize.interval\_drop}=2). The number of draws kept in the last iteration is \texttt{e(mcmc\_samplesize)} and the interval of that sample is \texttt{e(mcmc\_interval\_final)}. The explicit burn-in of each run is one interval; \opt{mcmcburnin()} (default 16 times the interval) is the burn-in of fixed-size samples, of simulation and of the goodness-of-fit simulations. The ESS of the target is that of the estimating equations (for a curved term, of the equations in the coefficients that are estimated), so that it is smaller than the ESS of an individual statistic as \texttt{estat mcmcdiag} reports it and is, as in \ergm{}, 64 and not several hundred; a failed convergence test raises the target and the base size (next paragraph). The sizes are capped at 32 times their starting values, a safeguard that \ergm{} does not have and that matters only in a fit that never converges.

\paragraph{The interplay of the options.} \opt{mcmcsamplesize(\#)} is the \emph{base} size of the adaptive sample, not the number of draws kept: the fit draws at least that many and more if the ESS target has not been reached. \opt{mcmleeffectivesize(\#)} sets the target ESS; a larger number gives a longer, more accurate sample, and \opt{mcmleeffectivesize(0)} switches the adaptive sample off, so that every iteration records exactly \opt{mcmcsamplesize()} draws (3000 if the option is not given) at \opt{mcmcinterval()} after \opt{mcmcburnin()}, the interval grows when the ESS of the statistics stays below the larger of 200 and 64 times the number of coefficients, the convergence test is a Hotelling test of the mean estimating equation, the covariance comes from one further sample at the final coefficients, and the bridge samples have 3000 draws; this is the fixed-size algorithm. The options \opt{mcmcinterval()} and \opt{mcmcburnin()} apply to both, in the adaptive sample as the starting values described above. Fits by contrastive divergence and stochastic approximation and fits with \opt{obsconstraints()} have no adaptive sample and use fixed-size samples whatever \opt{mcmleeffectivesize()} says (3000 draws unless \opt{mcmcsamplesize()} is given). A model whose chain does not mix can take very long: the sample is extended up to its caps, and 60 iterations can take many minutes. Lower \opt{mcmleiterations()} for a model that you suspect to be degenerate.

\paragraph{Stopping rules.} The option \opt{mcmletermination()} (R: \texttt{MCMLE.termination}), with the values \texttt{confidence}, \texttt{hummel}, \texttt{hotelling} and \texttt{none}, chooses how the fit stops. The default, \texttt{confidence}, is \ergm{}'s test: the fit stops when, with the confidence \opt{mcmleconfidence()} (R: \texttt{MCMLE.confidence}, default 0.99), the estimating equations at the updated coefficients lie inside a tolerance region, an ellipsoid whose size is 0.1 of the covariance of the equations (R: \texttt{MCMLE.MCMC.precision}), taking the autocorrelation of the chain into account. A fit that fails the test raises the target ESS by the factor min(critical value$/T^2$, 2) and the base sample size by its square root, as \ergm{} does, and iterates; this is why an adaptive fit can take several iterations. With \opt{mcmleeffectivesize(0)} the test is the earlier one, a Hotelling-type test that cannot distinguish the observed statistics from the mean of the sampled statistics at the confidence level. \texttt{hotelling} and \texttt{hummel} stop after the criterion has held in two successive iterations, with the sample size raised by a factor of four (\texttt{MCMLE.last.boost}) after the first; \texttt{none} runs all iterations. \opt{mcmlesteplength()} and \opt{mcmlestepmargin()} (R: \texttt{MCMLE.steplength}, \texttt{MCMLE.steplength.margin}) control the step length; the margin \texttt{.} switches the Hummel step length off. If the statistics of the sample do not vary at all the fit stops with R's error \emph{unconstrained MCMC sampling did not mix at all}, which is what happens for a degenerate model (Section~\ref{sec:kapferer}).

\paragraph{Standard errors.} The information matrix is estimated by the covariance of the simulated statistics of the last iteration's sample (as in \ergm{}; in a fixed-size fit, of one further sample drawn at the final coefficients, and in an adaptive fit that ran out of iterations before convergence, too). As in \ergm{}'s default, the reported covariance is the inverse information plus the Monte Carlo error of the estimate, $V^{-1}+V^{-1}(S/n)V^{-1}$, where $S$ is an autoregressive spectral estimate of the long-run covariance of the draws. Standard errors therefore grow when the chain mixes badly. A coefficient table whose standard errors are larger than those of a similar pseudolikelihood fit is usually not an error: the pseudolikelihood standard errors of a dyad-dependent model are too small.

\subsection{Log-likelihood, AIC and BIC}\label{sec:loglik}

\nwergm{} stores the log-likelihood as R's \texttt{logLik()}: exact for a dyad-independent model (where the pseudolikelihood is the likelihood), the pseudolikelihood after \opt{estimate(mple)} of a dyad-dependent model, and, after Monte Carlo maximum likelihood, the likelihood evaluated by \emph{bridge sampling} from the dyad-independent submodel (R: \texttt{ergm.bridge.dindstart.llk}). The log-likelihood ratio between the fitted model and the dyad-independent model that keeps only the dyad-independent terms is the integral of the expected statistic along a path between the two coefficient vectors; the integral is approximated by $B$ bridges, each with its own Markov chain. The options are \opt{bridgesteps()} (default 16, R: \texttt{bridge.nsteps}), \opt{bridgesamplesize()} (default the number of draws of the last iteration, as in \ergm{}, so a few hundred; with \opt{mcmleeffectivesize(0)}, \opt{mcmcsamplesize()}) and \opt{bridgeinterval()} (default the final thinning interval divided by the number of bridges); \opt{nologlik} skips the evaluation, in which case the log-likelihood, AIC and BIC are missing. The Monte Carlo standard error of the log-likelihood is \texttt{e(ll\_se)}; it is zero when the log-likelihood is exact. The null model is the model with all coefficients zero, whose log-likelihood is $N\ln\frac12$ with $N$ the number of informative dyads. The deviance is $-2$ times the log-likelihood, AIC and BIC use the number of estimated coefficients (offsets excluded) and, for the BIC, the number of informative dyads. The likelihood is not evaluated for a model with a curved term fitted by Monte Carlo maximum likelihood.

\subsection{Contrastive divergence}\label{sec:cd}

\opt{estimate(cd)} (R: \texttt{estimate="CD"}) is contrastive divergence. Starting from zero, each draw is a short Markov chain of \opt{cdnsteps()} proposals (R: \texttt{CD.nsteps}, default 8) started from the \emph{observed} network, and the coefficients are updated by an importance-sampling Newton step with the Hummel step length. The number of draws per coefficient vector is \opt{cdsamplesize()} (default 128 times the number of free coefficients, R: \texttt{CD.samplesize.per\_theta}), and \opt{mcmleiterations()} bounds the number of iterations (R: \texttt{CD.maxit}). Contrastive divergence is cheap and is meant to provide starting values for the Monte Carlo estimator; its estimates are not maximum likelihood estimates. R reports the covariance of the estimate as missing; because a Stata estimation result cannot hold missing values on the diagonal, \nwergm{} stores the covariance of the contrastive-divergence sample in \texttt{e(V)} and says so under the table.

\subsection{Stochastic approximation}\label{sec:sa}

\opt{mainmethod(sa)} (R: \texttt{main.method="Stochastic-Approximation"}) replaces the Monte Carlo maximum likelihood iteration by the stochastic approximation of Snijders~\cite{snijders2002}: phase~1 estimates a gain from short runs, phase~2 consists of four sub-phases of Robbins--Monro updates with decreasing gain, and phase~3 draws a sample at the final coefficients for a Newton step and the covariance. The controls are \opt{saburnin()} (R: \texttt{SA.burnin}, default 16384), \opt{sainterval()} (\texttt{SA.interval}, 1024), \opt{saphase3()} (\texttt{SA.phase3\_n}, 1000), \opt{sagain()} (\texttt{SA.initial\_gain}, 0.1) and \opt{sansub()} (\texttt{SA.nsubphases}, 4). Stochastic approximation is slower than Monte Carlo maximum likelihood on well-behaved models and is an alternative when the latter does not converge.

\subsection{Starting values and offsets}

\opt{init(numlist)} (R: \texttt{control.ergm(init=)}) gives one starting value per coefficient, in the order of the coefficient table; a period or \texttt{NA} leaves a value to be estimated. Values for coefficients that are fixed by \opt{offset()} replace the offset values; other values are the starting values of the Monte Carlo fit and are held fixed in the pseudolikelihood step that finds the missing ones. If every value is given, the pseudolikelihood step is skipped. With \opt{estimate(mple)} only offset values are used, except for curved terms, where \opt{init()} gives the starting point of the pseudolikelihood fit (one value for each parameter, the weight and decay of the curved term last).

\opt{offset(coefname \# [coefname \# \ldots])} (R: the \texttt{offset()} wrapper with \texttt{offset.coef}) fixes named coefficients at known values. For the pseudolikelihood, the contribution of the fixed term enters the logistic regression as an offset; for Monte Carlo maximum likelihood, the Newton step is restricted to the free coefficients, while the term still takes part in the simulation (fixing a coefficient does not remove its term from the model). An offset coefficient is reported at the given value with standard error zero and a zero row and column in \texttt{e(V)}. At least one coefficient must remain free, and offsets are not allowed for the weight and decay of a curved term.

\section{Reproducibility}\label{sec:seed}

Monte Carlo estimates depend on the random numbers. \opt{seed(\#)} sets the seed of the sampler; it is drawn from Stata's random-number generator right after \texttt{set seed}, so that the stream is deterministic. The same \texttt{nwergm} call with the same \opt{seed()}, the same data and the same plugin gives \emph{identical} coefficients in any session, whatever was loaded or fitted before. (An earlier release failed this: the starting network was built in the order in which the core returns the neighbors of a node, which differs once another network has been loaded; the neighbors are now read in sorted order, and a test fits the same model after loading other networks.) Section~\ref{sec:flo_seed} shows the guarantee at work.

What the seed does \emph{not} guarantee:
\begin{itemize}[nosep]
  \item A different operating system or plugin build, or a change of \opt{mcmcinterval()}, \opt{mcmcsamplesize()} or another control, gives a different (statistically equivalent) Monte Carlo fit.
  \item The stream is not the random-number stream of R. Results match \ergm{} only up to Monte Carlo error.
  \item Without \opt{seed()} the current Stata seed is used; two runs in the same session then differ.
\end{itemize}
Report the seed with the results, and judge the Monte Carlo error by refitting with other seeds (Section~\ref{sec:flo_seed} and Chapter~\ref{ch:validation} show how large the spread over seeds is).

\section{Missing dyads}\label{sec:missing}

A tie that is missing (a missing value in the edge matrix of the network, off the diagonal) is not treated as a non-tie. As in \ergm{}, \nwergm{} detects missing dyads automatically and
\begin{itemize}[nosep]
  \item drops them from the (pseudo)likelihood; the number of informative dyads \texttt{e(N)} excludes them, and for a dyad-independent model the pseudolikelihood estimate is the maximum likelihood estimate;
  \item for a dyad-dependent model fitted by Monte Carlo maximum likelihood, draws a second, constrained sample in which only the missing dyads are free and every observed dyad is held fixed (Handcock and Gile~\cite{handcock2010}); the estimating equations compare the two samples, and the information matrix is the difference of their covariances. The missing dyads of the starting network are imputed at random with the observed density, as \ergm{} does;
  \item reports the number of missing dyads in \texttt{e(nmissing)}.
\end{itemize}
A dyad-dependent model fitted by pseudolikelihood also imputes the missing dyads of the starting network at random (so the result depends on the seed, as in R). Missing dyads cannot yet be combined with \opt{freedyads()}, \opt{blockdiag()} or \opt{fixdensity}. \emph{Missing covariates} are not supported, as in \ergm{}: a node attribute with a missing value in a term (\opt{nodecov()}, \opt{nodematch()}, \ldots) and a dyadic covariate with missing dyads stop the command with return code 416. Drop or impute the missing values first. For two-mode terms only the attribute of the term's own mode is checked. An example is in Section~\ref{sec:missingex}.

\section{Constraints and observation processes}\label{sec:constraints}

\opt{constraints(\emph{formula})} restricts the set $\mathcal{Y}$ of networks on which the model is defined, exactly as \texttt{constraints=} does in \ergm{}. The formula is a list of constraint terms joined by \texttt{+} or \texttt{-}, with arguments written as in R, by position or by name. Two kinds of constraints exist. Constraints that \emph{fix some dyads} (\texttt{blocks}, \texttt{blockdiag}, \texttt{fixedas}, \texttt{fixallbut}, \texttt{observed}, \texttt{egocentric}) mark dyads that may not change; they are combined from left to right, a term with \texttt{+} keeping only the dyads that may vary according to it and to everything to its left, and a term with \texttt{-} adding the dyads that may vary according to it. Constraints that \emph{tie the dyads together} (\texttt{edges}, \texttt{degrees}, \texttt{odegrees}, \texttt{idegrees}, \texttt{b1degrees}, \texttt{b2degrees}, \texttt{degreedist}, \texttt{idegreedist}, \texttt{bd}) always restrict the sample space and may only carry a \texttt{+}. Constraints implied by others are dropped, as in R. The complete list is in Section~\ref{sec:conslist}.

Statistics that a constraint holds constant (for example \texttt{edges} under \texttt{constraints(edges)}, or \texttt{kstar()} under \texttt{constraints(degrees)}) are fixed at zero and ignored in the fit, with a note, as R does. A dyad-dependent constraint makes the sample space dyad-dependent, so that the default method is Monte Carlo maximum likelihood, started from zero. With \opt{estimate(mple)} the pseudolikelihood is formed over the free dyads.

\paragraph{Proposals.} The proposal is chosen from the constraints with R's rules. With no constraint, with dyad constraints and with \texttt{bd}, one dyad is toggled (tie/no-tie weighted or uniform, as above, with the counts of the proposal restricted to the free dyads). With \texttt{edges}, a tie is swapped with a non-tie. With \texttt{degrees}, or \texttt{idegrees} together with \texttt{odegrees}, or \texttt{b1degrees} together with \texttt{b2degrees}, tetrad moves (and, for directed networks, hexad moves) keep all degrees. With \texttt{odegrees}, \texttt{idegrees}, \texttt{b1degrees} or \texttt{b2degrees} alone, one end of a tie is moved. With \texttt{degreedist} or \texttt{idegreedist}, the ends of ties are exchanged between nodes whose degrees differ by one. With \texttt{dyadnoise}, a tie/no-tie proposal with the noise likelihood in the acceptance ratio is used. Other combinations stop with the message that R gives. The constraint options of the command are \opt{mcmcprop()} (R's \texttt{MCMC.prop} hints; \texttt{sparse} is the default, and \texttt{triadic} and \texttt{strat} are accepted but the tie/no-tie proposal is used, with a note, as R does when no proposal can use a hint) and \opt{mcmcmaxedges()} (R's \texttt{MCMC.maxedges}).

The older options \opt{freedyads(\emph{netname})} (the dyads that are ties in a network in memory may vary, every other dyad is fixed: R's \texttt{fixallbut()}), \opt{blockdiag(\emph{varname})} (only dyads within the same value of the variable may vary: R's \texttt{blockdiag()}) and \opt{fixdensity} (R's \texttt{constraints=\textasciitilde edges}) remain available; they can be given one at a time and cannot be combined with \opt{constraints()}.

\paragraph{Observation processes.} \opt{obsconstraints(\emph{formula})} (R: \texttt{obs.constraints}) describes how the network was \emph{observed}, for example the dyads that were not observed, as \texttt{fixallbut(\emph{unobserved})}. Two chains are run at every iteration, one under \opt{constraints()} and one under \opt{constraints()} and \opt{obsconstraints()} together, both started from the observed network. The coefficients solve the equation that sets the two expected statistics equal, and the information is the difference of the covariances of the two chains. This is the same estimator that handles missing dyads (Section~\ref{sec:missing}) with a general constraint in place of the set of missing dyads. It requires Monte Carlo maximum likelihood. The constraint \texttt{observed}, which frees exactly the missing dyads, is the default observation constraint for a network with missing dyads; in \opt{constraints()} on a network with missing dyads R itself stops with an internal error, and \nwergm{} stops with a plain message and points to \opt{obsconstraints(observed)}. On a network without missing dyads the empty constraint is dropped, as in R, and \nwergm{} prints a note.

\section{Curved terms}\label{sec:curved}

A curved term has a weight and a decay parameter (reported as, for example, \texttt{gwesp\_weight} and \texttt{gwesp\_decay}). The fit follows R: the pseudolikelihood fit is a trust-region optimization in the curved parameters, and the Monte Carlo fit takes a Newton step in the space of $\eta$ and projects it back onto the curve. A decay can be weakly identified on a small or sparse network, in which case the fit converges to a boundary solution (decay near 0, or very large). This is a property of the model, not a bug: R's \ergm{} converges to the same kind of solution on the same data and can fail to converge at all on a poorly identified network, and so can \nwergm{}, which then reports a non-converged fit rather than a plausible-looking wrong number (\texttt{e(converged)}). A large standard error of the decay is the usual sign; consider a longer burn-in, a larger sample, a different starting decay, or the fixed-decay form of the term. Fixed-decay terms (\opt{gwesp(0.5)}) are ordinary terms and do not have this problem. The option \opt{init()} gives starting values for the curved fit.
